\documentclass[conference]{IEEEtran}

% ── Encoding & Language ──────────────────────────────────────────────
\usepackage[utf8]{inputenc}
\usepackage[T1]{fontenc}
\usepackage[english]{babel}

% ── Math & Symbols ───────────────────────────────────────────────────
\usepackage{amsmath}
\usepackage{amssymb}

% ── Graphics & Figures ───────────────────────────────────────────────
\usepackage{graphicx}
\usepackage{float}

% ── Tables ───────────────────────────────────────────────────────────
\usepackage{booktabs}
\usepackage{tabularx}
\usepackage{array}

% ── Lists ────────────────────────────────────────────────────────────
\usepackage{enumitem}
\setlist[itemize]{leftmargin=1.5em}

% ── Code Blocks ──────────────────────────────────────────────────────
\usepackage{listings}
\usepackage{xcolor}
\definecolor{lightgray}{RGB}{245,245,245}
\lstset{
  basicstyle=\ttfamily\small,
  backgroundcolor=\color{lightgray},
  frame=single,
  breaklines=true
}

% ── Hyperlinks (load LAST) ───────────────────────────────────────────
\usepackage[
  colorlinks=true,
  linkcolor=black,
  urlcolor=black,
  citecolor=black,
  pdftitle={Project Document},
  pdfauthor={Author Name}
]{hyperref}

% ── Metadata ─────────────────────────────────────────────────────────
\title{Pose Imitation of Human Motion by a Simulated Humanoid Robot}

\makeatletter
\newcommand{\linebreakand}{%
  \end{@IEEEauthorhalign}
  \hfill\mbox{}\par
  \mbox{}\hfill\begin{@IEEEauthorhalign}
}
\makeatother

\author{
  \IEEEauthorblockN{1\textsuperscript{st} Given Name Surname}
  \IEEEauthorblockA{\textit{dept. name of organization (of Aff.)} \\
    \textit{name of organization (of Aff.)}\\
    City, Country \\
    email address}
  \and
  \IEEEauthorblockN{2\textsuperscript{nd} Given Name Surname}
  \IEEEauthorblockA{\textit{dept. name of organization (of Aff.)} \\
    \textit{name of organization (of Aff.)}\\
    City, Country \\
    email address}
  \and
  \IEEEauthorblockN{3\textsuperscript{rd} Given Name Surname}
  \IEEEauthorblockA{\textit{dept. name of organization (of Aff.)} \\
    \textit{name of organization (of Aff.)}\\
    City, Country \\
    email address}
  \linebreakand % <------------- \and with a line-break
  \IEEEauthorblockN{4\textsuperscript{th} Given Name Surname}
  \IEEEauthorblockA{\textit{dept. name of organization (of Aff.)} \\
    \textit{name of organization (of Aff.)}\\
    City, Country \\
    email address}
  \and
  \IEEEauthorblockN{5\textsuperscript{th} Given Name Surname}
  \IEEEauthorblockA{\textit{dept. name of organization (of Aff.)} \\
    \textit{name of organization (of Aff.)}\\
    City, Country \\
    email address}
}

% ─────────────────────────────────────────────────────────────────────
\begin{document}
 
\maketitle
 
% ─────────────────────────────────────────────────────────────────────
% ABSTRACT
% ─────────────────────────────────────────────────────────────────────
\begin{abstract}
This paper presents the design and implementation of a markerless visual teleoperation pipeline
that enables a simulated humanoid robot to imitate the live or recorded motion of a single human
using only a monocular RGB camera as input. The system processes a video stream through a
pose estimation module, maps the resulting keypoints to the robot's joint space via a retargeting
module, and drives the actuators of a simulated humanoid in Webots in near real time. The
pipeline targets an end-to-end latency of at most 150\,ms and a throughput of at least 20\,FPS
on commodity hardware. Imitation fidelity is evaluated using per-joint mean absolute error and
latency measurements. The work is conducted in the context of the CPSM 2026S course at
Technische Hochschule Mittelhessen under the supervision of M.Sc.\ Severin Stahl.
\end{abstract}
 
\begin{IEEEkeywords}
humanoid robot, pose estimation, motion imitation, Webots, teleoperation, MediaPipe, retargeting
\end{IEEEkeywords}
 
% ─────────────────────────────────────────────────────────────────────
% I. INTRODUCTION
% ─────────────────────────────────────────────────────────────────────
\section{Introduction}
Humanoid robots are a topic of strong scientific interest because they promise to operate in
environments designed for humans without costly infrastructure modifications. However, bipedal
locomotion and whole-body action require intelligent control to maintain balance. Recent research
increasingly leverages machine learning and human-motion imitation to address this challenge
\cite{ze2025twist}\cite{tao2021visual}.
 
This project explores markerless visual teleoperation of a humanoid robot inside the Webots
simulation environment. A single human is recorded by a tripod-mounted RGB camera; the captured
video is processed to extract 2D/3D human pose, mapped to the robot's joint space, and used to
drive the simulated humanoid's actuators in near real time.
 
\subsection{Purpose}
This document defines the project scope, system architecture, functional and non-functional
requirements, tooling, milestones, and acceptance criteria for the Pose Imitation project
(CPSM 2026S). It serves as the primary reference for all team members and the supervising
instructor throughout the development process.
 
\subsection{Scope}
The project considers a single human subject, fully visible in frame, recorded by a
tripod-mounted static RGB camera. Input may be either a live webcam stream or a pre-recorded
video file. Pose estimation and retargeting are implemented in Python, and a Webots Python
controller drives the robot's joints. Basic balance handling sufficient for upper-body and
slow lower-body imitation is included, along with quantitative and qualitative evaluation of
imitation fidelity and latency.
 
The following items are explicitly out of scope for the initial release: multi-person tracking,
dynamic or moving camera viewpoints, hardware deployment on a physical robot, highly dynamic
motions such as running or jumping, hand and finger-level dexterity, facial expression
imitation, and reinforcement-learning-based whole-body control. The latter may be considered
as future work.
 
% ─────────────────────────────────────────────────────────────────────
% II. PROBLEM STATEMENT
% ─────────────────────────────────────────────────────────────────────
\section{Problem Statement}
Enabling a humanoid robot to imitate human motion in real time is a complex challenge at the
intersection of computer vision, robotics, and control theory. Existing approaches often rely
on marker-based motion capture systems or depth sensors, which are expensive and limit
practical deployment. This project investigates whether a single monocular RGB camera,
combined with modern pose estimation libraries, can serve as a sufficient input for real-time
whole-body imitation in a simulated environment.
 
Markerless teleoperation removes the need for specialized hardware and enables intuitive,
low-cost robot control. Demonstrating this capability in simulation first allows rapid
iteration before any physical deployment. The project also serves as a practical investigation
into the limitations of monocular 3D pose estimation for robot control, including depth
ambiguity, self-occlusion, and latency constraints.
 
The following assumptions and constraints apply throughout the project. The human subject is
assumed to be fully visible in the camera frame at all times, and the camera is stationary.
The pipeline must be able to run on a modern laptop CPU, with GPU acceleration treated as
optional. The overall project timeline is fixed at ten weeks by the CPSM 2026S course schedule.
 
% ─────────────────────────────────────────────────────────────────────
% III. RELATED WORK
% ─────────────────────────────────────────────────────────────────────
\section{Related Work}
Ze et al.\ \cite{ze2025twist} present TWIST, a teleoperated whole-body imitation system that
drives a humanoid robot from human demonstrations, establishing the viability of
imitation-based control for complex manipulation and locomotion tasks.
 
Tao et al.\ \cite{tao2021visual} propose a visual perception method based on human pose
estimation for humanoid robot motion imitation, providing a direct methodological precedent
for the monocular RGB approach adopted in this project.
 
The present work differs from both references in that it targets a fully simulated environment
(Webots), uses exclusively open-source library-level pose estimation, and emphasizes
reproducibility and compatibility with low-cost hardware.
 
% ─────────────────────────────────────────────────────────────────────
% IV. SYSTEM OVERVIEW
% ─────────────────────────────────────────────────────────────────────
\section{System Overview}
 
\subsection{Vision and Stakeholders}
The goal is to build an end-to-end pipeline in which a simulated humanoid robot in Webots
imitates the live or recorded motion of a single human using only a monocular video stream
as input. Three stakeholder groups are involved: the student developers, who are responsible
for implementation, evaluation, and documentation; the supervising instructor M.Sc.\ Severin
Stahl, who provides requirements clarification, scientific guidance, and grading; and the
end user of the demo, who stands in front of the camera and observes the robot's imitation.
 
\subsection{System Architecture}
The pipeline consists of four sequential stages. The \textbf{Video Input Module} captures
timestamped frames from a webcam or video file using OpenCV \texttt{VideoCapture} at a
configurable resolution and frame rate. The \textbf{Pose Estimation Module} extracts per-frame
2D/3D body keypoints; candidate libraries are MediaPipe Pose (33 landmarks with 3D world
coordinates), MMPose, and Ultralytics YOLOv8-Pose. Temporal jitter is reduced using a
One-Euro filter or exponential smoothing. The \textbf{Retargeting Module} converts human
keypoints into joint angles compatible with the Webots humanoid by computing angles from
vector geometry between keypoints such as shoulder, elbow, hip, and knee, with optional
analytical or numerical IK for end-effector targets. Joint limits and rate limits are
enforced to ensure mechanically feasible commands. Finally, the \textbf{Webots Controller}
communicates with the simulation via the \texttt{controller} Python API, reads target joint
angles from the retargeting module, and issues \texttt{setPosition()} commands at each
control step. An optional simple balance assist using fixed feet or a PD-stabilized torso is
included for the initial milestones.
 
Regarding inter-process communication, the preferred option is to run the pose pipeline
inside the Webots controller process when performance allows. If latency becomes a concern,
the pipeline runs as a separate process streaming joint targets via a local socket, ZeroMQ,
or shared memory, with the choice determined by runtime benchmarks.
 
A \textbf{Logging and Evaluation} component runs in parallel, recording human keypoints,
target joint angles, achieved joint angles, and timestamps. Offline metrics include per-joint
MAE, end-to-end latency, dropped frames, and the simulation-to-real-time ratio.
 
% ─────────────────────────────────────────────────────────────────────
% V. TECHNOLOGIES AND TOOLS
% ─────────────────────────────────────────────────────────────────────
\section{Technologies and Tools}
The simulator is Webots R2023b or newer. The implementation language is Python 3.10 or newer.
Core libraries are \texttt{opencv-python} for camera capture and video I/O,
\texttt{mediapipe} or \texttt{ultralytics} for pose estimation, \texttt{numpy} and
\texttt{scipy} for numerical computation and IK, \texttt{pyyaml} for configuration parsing,
and \texttt{matplotlib} for metric plots and analysis. The \texttt{controller} Python module
shipped with Webots handles simulation communication. Code quality is enforced via
\texttt{ruff} for linting and \texttt{black} for formatting, with \texttt{pytest} covering
unit tests for retargeting mathematics. \texttt{pyzmq} is included as an optional dependency
for inter-process messaging.
 
% ─────────────────────────────────────────────────────────────────────
% VI. REQUIREMENTS
% ─────────────────────────────────────────────────────────────────────
\section{Requirements}
 
\subsection{Functional Requirements}
The system shall accept a live webcam feed and a recorded video file as input sources,
selectable via a single YAML or TOML configuration file (FR-1, FR-8). From each input frame,
the system shall extract human body keypoints (FR-2) and map them to the simulated humanoid's
joint space (FR-3). The Webots controller shall command the humanoid's motors to follow the
mapped joint trajectory at each simulation step (FR-4), clipping all commands to the robot's
mechanical limits (FR-5) and applying temporal smoothing to prevent unsafe oscillations
(FR-6). All human keypoints and robot joint states shall be logged with timestamps to disk
(FR-7). The complete system shall be launchable via a single command, \texttt{python run.py}
(FR-9).
 
\subsection{Non-Functional Requirements}
End-to-end latency from camera frame to joint command shall not exceed 150\,ms as a goal and
300\,ms as an acceptable upper bound (NFR-1). The pose estimation pipeline shall sustain at
least 20\,FPS on a modern laptop CPU or GPU (NFR-2). Upper-body imitation fidelity, measured
as per-joint MAE on representative test motions, shall remain at or below 10$^{\circ}$
(NFR-3). The robot shall not fall during standing upper-body imitation in the baseline scenario
(NFR-4). Reproducibility is ensured via fixed random seeds and pinned dependency versions
(NFR-5). The codebase shall include type hints, pass \texttt{ruff} and \texttt{black} checks,
and provide unit tests for retargeting mathematics (NFR-6). The pipeline shall run on macOS
and Linux with Webots R2023b or newer (NFR-7).
 
% ─────────────────────────────────────────────────────────────────────
% VII. USER STORIES
% ─────────────────────────────────────────────────────────────────────
\section{User Stories}
\begin{itemize}
  \item \textit{As a user}, I can stand in front of a webcam and see the simulated humanoid mirror my upper-body motions in Webots in near real time.
  \item \textit{As a developer}, I can replay a recorded video and have the robot imitate the same motion deterministically for evaluation purposes.
  \item \textit{As a researcher}, I can log joint trajectories for both the human and the robot and compute imitation-fidelity metrics offline.
  \item \textit{As a user}, the robot should not fall over during normal standing and upper-body imitation.
\end{itemize}
 
% ─────────────────────────────────────────────────────────────────────
% VIII. GOALS AND SUCCESS METRICS
% ─────────────────────────────────────────────────────────────────────
\section{Goals and Success Metrics}
The project aims to deliver a functional end-to-end imitation pipeline within the ten-week
course timeline, achieve real-time upper-body imitation without robot falls in the baseline
scenario, provide reproducible quantitative evaluation of latency and imitation fidelity, and
produce documentation sufficient for an independent developer to reproduce the results.
 
The project is considered complete when all of the following acceptance criteria are met:
\begin{bracketlist}
  \item Running a single command launches Webots and the Python pipeline together.
  \item With a live webcam, the simulated humanoid visibly imitates the user's upper-body motion in real time without falling.
  \item Replaying a recorded reference video reproduces equivalent robot motion deterministically.
  \item Quantitative metrics (latency, per-joint MAE) are reported in \texttt{docs/} along with plots.
  \item The repository contains documentation sufficient for a new developer to reproduce the results.
\end{bracketlist}
 
% ─────────────────────────────────────────────────────────────────────
% IX. PROJECT TIMELINE
% ─────────────────────────────────────────────────────────────────────
\section{Project Timeline}
 
\begin{table}[H]
  \centering
  \caption{Project Milestones}
  \begin{tabularx}{\columnwidth}{lXl}
    \toprule
    \textbf{ID} & \textbf{Milestone / Deliverable} & \textbf{Target} \\
    \midrule
    M1 & Environment setup: Webots world loads; controller stub commands one motor. & Week 1--2 \\
    M2 & Pose estimation prototype: keypoints extracted from webcam/video. & Week 3 \\
    M3 & Retargeting v1: shoulder/elbow angles drive the robot's arms. & Week 4--5 \\
    M4 & Real-time integration: live webcam $\rightarrow$ live robot imitation. & Week 6 \\
    M5 & Lower body and stability: hip/knee/ankle imitation with balance assist. & Week 7--8 \\
    M6 & Evaluation and logging: metrics, plots, recorded demo videos. & Week 9 \\
    M7 & Final report and cleanup: documentation and reproducible demo. & Week 10 \\
    \bottomrule
  \end{tabularx}
\end{table}
 
% ─────────────────────────────────────────────────────────────────────
% X. RISKS AND MITIGATIONS
% ─────────────────────────────────────────────────────────────────────
\section{Risks and Mitigations}
Several technical risks have been identified. Monocular 3D pose estimation is inherently
ambiguous due to depth uncertainty and self-occlusion, which may lead to inaccurate
retargeting. This is mitigated by using MediaPipe World Landmarks and enforcing joint limits.
Direct joint mapping without balance control may cause the robot to fall; the mitigation is
to start with feet fixed and add PD torso stabilization before enabling lower-body joints.
If end-to-end latency exceeds acceptable bounds, the pose estimation pipeline will be moved
to a separate process, input frames will be downscaled, and a GPU-accelerated model will be
used. Differences between the Webots humanoid joint structure and the human skeleton will be
addressed by building and validating an explicit, documented per-joint mapping table. Finally,
library and Webots version drift is mitigated by pinning all dependency versions in
\texttt{requirements.txt} and documenting the required Webots release explicitly.
 
% ─────────────────────────────────────────────────────────────────────
% XI. OPEN QUESTIONS
% ─────────────────────────────────────────────────────────────────────
\section{Open Questions}
The following questions remain open and are to be resolved with the supervising instructor
prior to the start of implementation. First, regarding the robot model: should the project
use a built-in Webots humanoid such as NAO or Atlas, or a custom PROTO model? Second, are
there required reference motions or poses against which imitation fidelity must be evaluated?
Third, is real-time imitation a hard grading requirement, or is offline replay of a recorded
video acceptable? Fourth, will a GPU be available for pose estimation, or must the pipeline
run on CPU only? Fifth, is a stretch goal of multi-camera input of interest to the supervisor?
 
% ─────────────────────────────────────────────────────────────────────
% XII. CONCLUSION
% ─────────────────────────────────────────────────────────────────────
\section{Conclusion}
This document has defined the scope, architecture, requirements, and timeline for the Pose
Imitation project. The proposed pipeline addresses the challenge of markerless humanoid robot
teleoperation using only a commodity RGB camera, offering a low-cost alternative to
marker-based motion capture. If the outlined milestones are met, the project will deliver a
reproducible, quantitatively evaluated demonstration of real-time human-to-robot motion
imitation in Webots. The open questions identified in Section~XI are to be resolved with the
supervising instructor before implementation begins.
 
% ─────────────────────────────────────────────────────────────────────
% REFERENCES
% ─────────────────────────────────────────────────────────────────────
\begin{thebibliography}{9}
 
\bibitem{ze2025twist}
Z.~Ze et al., ``TWIST: Teleoperated Whole-Body Imitation System,''
in \textit{Proc.\ Conference on Robot Learning (CoRL)}, 2025.
[Online]. Available: \url{https://doi.org/10.48550/arXiv.2505.02833}
 
\bibitem{tao2021visual}
Y.~Tao et al., ``Visual Perception Method Based on Human Pose Estimation for Humanoid
Robot Imitating Human Motions,'' in \textit{Proc.\ 2nd Int.\ Conf.\ Control, Robotics and
Intelligent System (CCRIS)}, 2021, pp.~54--61.
[Online]. Available: \url{https://doi.org/10.1145/3483845.3483867}
 
\end{thebibliography}
 
% ─────────────────────────────────────────────────────────────────────
\end{document}